\documentclass[10pt,a4paper]{amsart}
\usepackage[margin=19mm]{geometry}
\usepackage{amsmath,amssymb,mathtools}
\usepackage[scr=rsfs,cal=euler]{mathalfa}
\usepackage{xfrac}
\usepackage[unicode,hidelinks,bookmarks=false]{hyperref}
\newcommand{\ie}{i.e.\ }
\newcommand{\cf}{cf.\ }
\newcommand{\resp}{respectively\ }
\newcommand{\Subsection}{\subsection}
\providecommand{\nobreakdash}{\nobreak\hbox{-}}
\setlength{\emergencystretch}{2em}
\multlinegap0pt
\usepackage{bm}
\usepackage{bbm}
\usepackage{dsfont}
\usepackage{xspace}
\usepackage{cancel}
\usepackage{mathtools}
\usepackage{amsthm}
\makeatletter
\@ifundefined{theorem}{\newtheorem{theorem}{Theorem}}{}
\makeatother
\DeclareMathOperator{\Tr}{Tr} 
\newcommand{\Hd}{\mathcal{H}}
\newcommand{\R}{\mathbb{R}}
\newcommand{\contvf}{C^1_c(\R^n; \R^n)}
\newcommand{\dd}{\: d}
\newcommand{\divv}{\textnormal{div} \;}
\newcommand{\ind}{\mathbbm{1}}
\newcommand{\pom}{{\partial \Omega}}
\title[Optimal constant for the trace inequality in $BV$ for domain]{Optimal constant for the trace inequality in $BV$ for domains with corners}
\author{Riccardo Cristoferi, Devin van der Gulik}
\date{Source: arXiv:2604.13770v1 (CC BY 4.0)}
\begin{document}
\maketitle

\noindent\emph{Source and licence.} Optimal constant for the trace inequality in $BV$ for domains with corners. Version arXiv:2604.13770v1 (CC BY 4.0), distributed under the Creative Commons Attribution 4.0 International licence, as recorded in the arXiv licence metadata for this exact version and shown on the abstract page https://arxiv.org/abs/2604.13770v1. Full rights notice: licenses/arxiv-2604.13770-CC-BY-4.0.txt.\par\smallskip

\noindent\textit{Editorial scope.} Counted unit: the Theorem whose source label is \texttt{thm:Giusti} in the article (source0.tex lines 648--650, source label \texttt{thm:Giusti}), delivered together with the article's own complete original proof of it (source0.tex lines 652--692, one \texttt{proof} environment of 41 lines). No mathematics is written, repaired or re-derived: statement and proof are the article's own text line for line. Required context retained uncounted: thm:Giorgo (lines 523--533). Every definition, display and label of those lines is retained unchanged. Omitted from this package and not redistributed: 1--522, 534--647, 651--651, 693--1016. Nothing inside a retained excerpt is omitted or altered: every retained body is delivered exactly as the article's lines are written, so no omission receipt of any kind accompanies this package. Citations: the retained excerpts carry no citation command at all, so no bibliography entry of the article is transcribed and no reference is redistributed. Reference adaptation: none was needed and none was made; every reference inside the delivered unit resolves inside it, so no unresolved reference or citation is produced. Printed slips kept literal: no printed word, symbol, hypothesis or number of the counted statement or of its proof was altered. Layout and class: the article's own class is \texttt{amsart} with packages including lipsum, amsfonts, graphicx, epstopdf, algorithmic, inputenc, mathtools, amssymb, hyperref, textcomp, babel, verbatim, todonotes, bbm; the wrapper is the lane's pinned \texttt{amsart} editorial wrapper at 10pt on A4, which restates the theorem-like environments the retained text needs and adds the article's own macro definitions that the retained text needs (\texttt{\textbackslash Hd}, \texttt{\textbackslash R}, \texttt{\textbackslash Tr}, \texttt{\textbackslash contvf}, \texttt{\textbackslash dd}, \texttt{\textbackslash divv}, \texttt{\textbackslash ind}, \texttt{\textbackslash pom}), with extra wrapper packages (bm, bbm, dsfont, xspace, cancel, mathtools). The delivered numbering is the unit's own. Assets: no retained span contains a figure environment, a graphics-inclusion command, a PDF-page inclusion, a table or a TikZ picture; no image file is copied into this package, the package declares no asset dependency and no image file is redistributed with it. Licence: version arXiv:2604.13770v1 of this article is distributed under the Creative Commons Attribution 4.0 International licence, as recorded in the arXiv licence metadata for this exact version and shown on the abstract page https://arxiv.org/abs/2604.13770v1. A full rights notice is delivered at licenses/arxiv-2604.13770-CC-BY-4.0.txt. Characters: every retained excerpt is pure ASCII, so the delivered engine, XeLaTeX with its default Latin Modern text font, reports no missing character.
	\begin{theorem}[De Giorgi's structure theorem] \label{thm:Giorgo}
		If a set $E$ of locally finite perimeter in $\R^n$, then the Gauss-Green measure $\mu_E$ of $E$ satisfies
		\[
		\mu_E = \nu_E \Hd^{n-1} \llcorner \partial^\ast E, \quad |\mu_E| = \Hd^{n-1} \llcorner \partial^\ast E.
		\]
		Moreover, the generalized Gauss-Green formula holds true:
		\[
		\int_E \divv T \coloneqq \int_{\partial^\ast E} T \, \nu_E \, \dd \Hd^{n-1},
		\]
		for all $T \in C^1_c(\R^n;\R^n)$.
	\end{theorem}

	\begin{theorem} \label{thm:Giusti}
		Let $\pom$ be $C^1$ in a neighbourhood of $x_0 \in \pom$. Then, $q_{\pom}(x_0) = 1$.
	\end{theorem}

	\begin{proof}
		We can assume that $x_0=0$ and that $\pom$ can be represented as the graph of a $C^1$ function $f(x')$ for $x'=(x_1, \cdots, x_{n-1})$ such that $f(0)=0$ and $\nabla f(0) = 0$ and
		\[
		x_n > f(x') \quad \textnormal{ for } x \in \Omega \cap B(0,r) 
		\]
		Let $E \subset \Omega$ and let $\pi_E$ be the projection of  $\partial^\ast E \cap \pom$ on the hyperplane $x_n=0$. Then we have 
		\begin{equation*}
			\int_{\pom} \Tr(\ind_E) \: \dd \Hd^{n-1} = \Hd^{n-1}(\partial^\ast E \cap \pom) = \int_{\pi_E} \sqrt{1+|\nabla f(x')|^2} \dd x'
		\end{equation*}
		Furthermore setting $M_r \coloneqq \sup \{|\nabla f|(x), |x'| < r \}$ we get $M_r \rightarrow 0$ as $r \rightarrow 0$ and 
		\begin{equation} \label{eq:C1proof1}
			\int_\pom \Tr(\ind_E) \dd \Hd^{n-1} \leq
			\Bigl(\sqrt{1+M_r^2}\Bigr)\Hd^{n-1}(\pi_E)
			\leq
			(1+M_r)\Hd^{n-1}(\pi_E)
		\end{equation} 
		Now, let $T  \equiv e_n = (0, \cdots, 0, 1) \in \contvf$ be a constant vector field. Then by Theorem \ref{thm:Giorgo} we have
		\[
		0 = \int_E \divv T = \int_{\partial^\ast E \cap \Omega} T \nu_E \Hd^{n-1} + \int_{\partial^\ast E \cap \partial \Omega} T \nu_E \Hd^{n-1},  
		\]
		therefore  we get 
		\[
		- \int_{\partial^\ast E \cap \partial \Omega} T \cdot \nu_E \Hd^{n-1} = \int_{\partial^\ast E \cap \Omega} T \cdot \nu_E \Hd^{n-1} \leq \mu_E(\Omega) = P(E;\Omega). 
		\]
		On the other hand, 
		\begin{align*}
			-\int_{\partial^\ast E \cap \partial \Omega} T \cdot \nu_E \Hd^{n-1} &= -\int_{\pi_E} T \cdot \nu_E |\text{det } Df(x)| \dd x \\
			&= -\int_{\pi_E} T \cdot \frac{(\nabla f, -1)}{\sqrt{1+|\nabla f|^2}} \sqrt{1+|\nabla f|^2} \dd x \\
			&= - \int_{\pi_E} -1 \dd x \\
			&= \Hd^{n-1}(\pi_E).
		\end{align*}
		As such we find
		\begin{equation} \label{eq:C1proof2}
			P(E;\Omega) \geq \Hd^{n-1}(\pi_E)
		\end{equation}
		Combining \eqref{eq:C1proof1} and \eqref{eq:C1proof2} we get
		\begin{equation}
			\frac{1}{P(E;\Omega)} \int_\pom \Tr(\ind_E)(x) \dd\Hd^{n-1} \leq (1+M_r)
		\end{equation}
		If we now let $r \rightarrow 0$ the result follows.
	\end{proof}

\end{document}
